\chapter{Vettori aleatori e vettori gaussiani}

Nel Capitolo~14 abbiamo studiato coppie e $m$-ple di variabili continue. Quando le variabili diventano molte (i
campioni di un segnale in istanti successivi, le tensioni su $n$ antenne, i pixel di un'immagine) scrivere ogni media e
ogni covarianza con i suoi indici diventa pesantissimo. L'algebra lineare offre un modo compatto per farlo: si
raccolgono le variabili in un \textbf{vettore}, le loro medie in un altro vettore e le covarianze di tutte le coppie in
una \textbf{matrice}, e le formule tornano corte come quelle di una sola variabile. Con questa notazione studieremo la
famiglia più importante di tutte, i \textbf{vettori gaussiani}: $n$ variabili gaussiane legate fra loro, descritte da una
sola formula e con una proprietà preziosa: restano gaussiane quando passano attraverso un sistema lineare.

Come nelle slide, i vettori si indicano con una sottolineatura ($\vett{X}$, $\vett{x}$, $\vett{\mu}$) e le matrici
con due ($\matr{A}$, $\matr{C}$).

% ============================================================
\section{Vettori e matrici: il minimo indispensabile}
% ============================================================

Per leggere questo capitolo bastano poche operazioni fra vettori e matrici. Le raccogliamo qui, partendo da zero; chi
le conosce già può passare direttamente ai vettori aleatori.

\begin{richiamo}{vettori, matrici e trasposti}
Un \term{vettore} di $n$ numeri si scrive come una \term{colonna}:
\[ \vett{x} = \begin{bmatrix} x_1 \\ x_2 \\ \vdots \\ x_n \end{bmatrix} \qquad (n \times 1:\ n \text{ righe e una
   colonna}) . \]
Il suo \term{trasposto} $\vett{x}^T = [x_1, x_2, \ldots, x_n]$ contiene gli stessi numeri disposti in riga
($1 \times n$). Per risparmiare spazio, un vettore colonna si scrive spesso come il trasposto di una riga:
$\vett{x} = [x_1, x_2, \ldots, x_n]^T$.

Una \term{matrice} $m \times n$ è una tabella di numeri con $m$ righe e $n$ colonne; l'elemento $a_{ij}$ sta nella riga
$i$ e nella colonna $j$. La \term{trasposta} $\matr{A}^T$ si ottiene scambiando le righe con le colonne: l'elemento di
posto $(i, j)$ di $\matr{A}^T$ è $a_{ji}$. Per esempio
\[ \matr{A} = \begin{bmatrix} 1 & 2 & 3 \\ 4 & 5 & 6 \end{bmatrix} \ (2 \times 3) , \qquad
   \matr{A}^T = \begin{bmatrix} 1 & 4 \\ 2 & 5 \\ 3 & 6 \end{bmatrix} \ (3 \times 2) . \]
Una matrice quadrata ($n \times n$) è \term{simmetrica} se $\matr{A}^T = \matr{A}$, cioè se $a_{ij} = a_{ji}$: è uguale
alla sua immagine allo specchio rispetto alla \term{diagonale principale}, quella che va dall'angolo in alto a sinistra
all'angolo in basso a destra. È \term{diagonale} se tutti gli elementi fuori dalla diagonale principale sono nulli.
\end{richiamo}

\begin{richiamo}{i prodotti «riga per colonna»}
Tutti i prodotti fra matrici e vettori seguono una sola regola, \textbf{riga per colonna}: si prende una riga del primo
fattore e una colonna del secondo, si moltiplicano gli elementi a due a due e si sommano i risultati.
\begin{itemize}
  \item \textbf{Riga per colonna} (prodotto scalare): $\vett{x}^T \vett{y} = x_1 y_1 + x_2 y_2 + \cdots + x_n y_n$. Il
        risultato è un \textbf{numero} ($1 \times n$ per $n \times 1$ dà $1 \times 1$). Per esempio
        $\vett{x}^T \vett{x} = \sum_i x_i^2$ è il quadrato della lunghezza del vettore (per Pitagora).
  \item \textbf{Matrice per vettore}: $\matr{A}\vett{x}$ è il vettore che ha per elemento $i$-esimo il prodotto della
        riga $i$ di $\matr{A}$ per $\vett{x}$: $(\matr{A}\vett{x})_i = \sum_j a_{ij} x_j$. Ogni elemento del risultato è
        una \textbf{combinazione lineare} delle $x_j$ (figura qui sotto).
  \item \textbf{Matrice per matrice}: l'elemento $(i, j)$ di $\matr{A}\matr{B}$ è la riga $i$ di $\matr{A}$ per la
        colonna $j$ di $\matr{B}$. Le dimensioni «interne» devono coincidere e spariscono:
        $(m \times n)(n \times p) = m \times p$. Attenzione: in generale $\matr{A}\matr{B} \neq \matr{B}\matr{A}$,
        l'ordine dei fattori conta.
  \item \textbf{Colonna per riga}: $\vett{x}\vett{y}^T$ ($n \times 1$ per $1 \times n$) è una \textbf{matrice}
        $n \times n$, con elemento $(i, j)$ uguale a $x_i y_j$. In particolare $\vett{x}\vett{x}^T$ contiene
        \textbf{tutti i prodotti a due a due} $x_i x_j$:
        \[ \vett{x}\vett{x}^T = \begin{bmatrix} x_1 \\ x_2 \end{bmatrix} [x_1, x_2] = \begin{bmatrix} x_1^2 & x_1 x_2
           \\ x_2 x_1 & x_2^2 \end{bmatrix} . \]
\end{itemize}
\begin{center}
\begin{tikzpicture}[ampersand replacement=\&, every left delimiter/.style={xshift=0.9ex},
  every right delimiter/.style={xshift=-0.9ex}]
  \matrix (A) [matrix of math nodes, left delimiter={[}, right delimiter={]},
    nodes={minimum width=2.3em, minimum height=1.6em, font=\small}] {
    a_{11} \& a_{12} \& a_{13} \\
    a_{21} \& a_{22} \& a_{23} \\
  };
  \matrix (x) [matrix of math nodes, left delimiter={[}, right delimiter={]}, right=0.5cm of A,
    nodes={minimum width=1.8em, minimum height=1.6em, font=\small}] {
    x_1 \\ x_2 \\ x_3 \\
  };
  \node[right=0.35cm of x] (uguale) {$=$};
  \matrix (y) [matrix of math nodes, left delimiter={[}, right delimiter={]}, right=0.35cm of uguale,
    nodes={minimum height=1.6em, font=\small}] {
    a_{11} x_1 + a_{12} x_2 + a_{13} x_3 \\
    a_{21} x_1 + a_{22} x_2 + a_{23} x_3 \\
  };
  \begin{scope}[on background layer]
    \node[fill=evid, rounded corners=2pt, inner sep=0pt, fit=(A-2-1)(A-2-3)] {};
    \node[fill=evid, rounded corners=2pt, inner sep=0pt, fit=(x-1-1)(x-3-1)] {};
    \node[fill=evid, rounded corners=2pt, inner sep=0pt, fit=(y-2-1)] {};
  \end{scope}
  \node[font=\scriptsize, anchor=north] at ([yshift=-2mm]A.south |- x.south) {$\matr{A}$ \ ($2 \times 3$)};
  \node[font=\scriptsize, anchor=north] at ([yshift=-2mm]x.south) {$\vett{x}$ \ ($3 \times 1$)};
  \node[font=\scriptsize, anchor=north] at ([yshift=-2mm]y.south |- x.south) {$\matr{A}\vett{x}$ \ ($2 \times 1$)};
  \node[font=\small, anchor=west, text width=4.6cm] at ([xshift=0.6cm]y.east)
    {La seconda riga di $\matr{A}$ per la colonna $\vett{x}$ dà il secondo elemento del risultato.};
\end{tikzpicture}
\end{center}
\textbf{Il trasposto di un prodotto} si ottiene trasponendo i fattori e \textbf{invertendone l'ordine}:
$(\matr{A}\matr{B})^T = \matr{B}^T \matr{A}^T$. Il perché: l'elemento $(i, j)$ di $(\matr{A}\matr{B})^T$ è l'elemento
$(j, i)$ di $\matr{A}\matr{B}$, cioè la riga $j$ di $\matr{A}$ per la colonna $i$ di $\matr{B}$; ma la colonna $i$ di
$\matr{B}$ è la riga $i$ di $\matr{B}^T$, e la riga $j$ di $\matr{A}$ è la colonna $j$ di $\matr{A}^T$. Per esempio
$(\matr{A}\vett{x})^T = \vett{x}^T \matr{A}^T$. Infine, un numero è uguale al suo trasposto, quindi
$\vett{x}^T \vett{y} = (\vett{x}^T \vett{y})^T = \vett{y}^T \vett{x}$.
\end{richiamo}

\begin{richiamo}{identità, inversa, determinante e forme quadratiche}
\textbf{Identità e inversa.} La matrice \term{identità} $\matr{I}$ ha 1 sulla diagonale principale e 0 altrove: per ogni
vettore $\matr{I}\vett{x} = \vett{x}$, come moltiplicare un numero per 1. L'\term{inversa} di una matrice quadrata
$\matr{A}$ è la matrice $\matr{A}^{-1}$ tale che $\matr{A}^{-1}\matr{A} = \matr{A}\matr{A}^{-1} = \matr{I}$: disfa quello
che fa $\matr{A}$, come $\frac{1}{a}$ disfa la moltiplicazione per $a$. Se $\vett{z} = \matr{A}\vett{x}$, allora
$\vett{x} = \matr{A}^{-1}\vett{z}$. Per un prodotto l'ordine si inverte, come per il trasposto:
$(\matr{A}\matr{B})^{-1} = \matr{B}^{-1}\matr{A}^{-1}$ (per disfare due operazioni si disfa prima l'ultima, come quando
ci si toglie prima le scarpe e poi i calzini). Inoltre $(\matr{A}^T)^{-1} = (\matr{A}^{-1})^T$.

\textbf{Determinante.} A ogni matrice quadrata si associa un numero, il \term{determinante}. Per una matrice $2 \times 2$:
\[ \det \begin{bmatrix} a & b \\ c & d \end{bmatrix} = ad - bc , \qquad
   \begin{bmatrix} a & b \\ c & d \end{bmatrix}^{-1} = \frac{1}{ad - bc} \begin{bmatrix} d & -b \\ -c & a \end{bmatrix} \]
(per verificare la formula dell'inversa basta moltiplicare: si ottiene $\matr{I}$). Il significato geometrico:
$\abs*{\det \matr{A}}$ è il \textbf{fattore per cui $\matr{A}$ moltiplica le aree} (i volumi, in $n$ dimensioni). La
matrice è invertibile esattamente quando $\det \matr{A} \neq 0$: se il determinante è nullo, $\matr{A}$ schiaccia il
piano su una retta e non si può più tornare indietro. Valgono $\det(\matr{A}\matr{B}) = \det\matr{A} \cdot \det\matr{B}$
(le due dilatazioni si susseguono) e $\det \matr{A}^T = \det \matr{A}$. Per una matrice diagonale tutto è facile: il
determinante è il prodotto degli elementi della diagonale, e l'inversa si ottiene invertendoli uno per uno.
\begin{center}
\begin{tikzpicture}[scale=0.95]
  \draw[-{Stealth}, mspigray] (-0.3,0) -- (2,0) node[below, font=\scriptsize, black] {$x_1$};
  \draw[-{Stealth}, mspigray] (0,-0.3) -- (0,2.1) node[left, font=\scriptsize, black] {$x_2$};
  \fill[evid] (0,0) rectangle (1,1); \draw[mspiric, line width=0.8pt] (0,0) rectangle (1,1);
  \node[font=\scriptsize] at (0.5,0.5) {area 1};
  \node[font=\scriptsize, below] at (1,0) {1}; \node[font=\scriptsize, left] at (0,1) {1};
  \draw[-{Stealth}, line width=0.8pt] (2.4,1) -- node[above, font=\small] {$\matr{A}$} (3.6,1);
  \begin{scope}[xshift=4.3cm]
    \draw[-{Stealth}, mspigray] (-0.3,0) -- (3.4,0) node[below, font=\scriptsize, black] {$z_1$};
    \draw[-{Stealth}, mspigray] (0,-0.3) -- (0,2.3) node[left, font=\scriptsize, black] {$z_2$};
    \fill[evid] (0,0) -- (2,0.5) -- (3,2) -- (1,1.5) -- cycle;
    \draw[mspiric, line width=0.8pt] (0,0) -- (2,0.5) -- (3,2) -- (1,1.5) -- cycle;
    \node[font=\scriptsize] at (1.5,1) {area $2{,}5$};
    \draw[-{Stealth}, mspiA, line width=0.9pt] (0,0) -- (2,0.5) node[below right, font=\scriptsize] {$(2;\ 0{,}5)$};
    \draw[-{Stealth}, mspiB, line width=0.9pt] (0,0) -- (1,1.5) node[above left, font=\scriptsize] {$(1;\ 1{,}5)$};
  \end{scope}
  \node[font=\small, anchor=west, text width=7.4cm] at (8.4,1)
    {Con $\matr{A} = \begin{bmatrix} 2 & 1 \\ 0{,}5 & 1{,}5 \end{bmatrix}$ il quadrato di lato 1 diventa il
     parallelogramma che ha per lati le colonne di $\matr{A}$; la sua area è
     $\det\matr{A} = 2 \cdot 1{,}5 - 1 \cdot 0{,}5 = 2{,}5$.};
\end{tikzpicture}
\end{center}
\textbf{Forme quadratiche.} Data una matrice quadrata $\matr{C}$, il numero
$\vett{q}^T \matr{C} \vett{q} = \sum_i \sum_j q_i\, c_{ij}\, q_j$ si chiama \term{forma quadratica}: è un polinomio di
secondo grado nelle $q_i$. Se $\matr{C}$ è diagonale, $\vett{q}^T \matr{C} \vett{q} = c_{11} q_1^2 + c_{22} q_2^2 + \cdots
+ c_{nn} q_n^2$: una somma di quadrati pesati.
\end{richiamo}

% ============================================================
\section{Vettori aleatori: medie, correlazioni e covarianze}
% ============================================================

\dfn{Vettore aleatorio}{
  Un \term{vettore aleatorio} $n$-dimensionale è la colonna di $n$ variabili aleatorie
  \[ \vett{X} = [X_1, X_2, \ldots, X_n]^T . \]
  La sua CDF e la sua PDF sono la CDF e la PDF congiunte delle $n$ variabili, scritte in forma vettoriale:
  $F_{\vett{X}}(\vett{x}) = P(X_1 \leq x_1, \ldots, X_n \leq x_n)$ e
  $f_{\vett{X}}(\vett{x}) = f_{X_1 \cdots X_n}(x_1, \ldots, x_n)$, con $\vett{x} = [x_1, \ldots, x_n]^T$.
}

Niente di nuovo, quindi: $f_{\vett{X}}(\vett{x})$ è la densità congiunta del Capitolo~14 con $n$ variabili, scritta in
modo compatto. La novità è la caratterizzazione sintetica. La media di un vettore (o di una matrice) di variabili
aleatorie si calcola \textbf{elemento per elemento}.

\dfn{Vettore delle medie}{
  Il \term{vettore valore medio} di $\vett{X}$ è il vettore colonna dei valori attesi delle $n$ variabili:
  \[ \vett{\mu}_{\vett{X}} \stackrel{\triangle}{=} E[\vett{X}] = [\mu_{X_1}, \mu_{X_2}, \ldots, \mu_{X_n}]^T . \]
}

Le medie da sole non dicono come le variabili sono legate fra loro. Per questo si considerano anche la correlazione e
la covarianza di tutte le coppie $(X_i, X_j)$ che si possono estrarre dal vettore, raccolte in una tabella $n \times n$,
cioè in una matrice. Il modo per costruirla in un colpo solo è il prodotto colonna per riga $\vett{X}\vett{X}^T$, che
contiene tutti i prodotti $X_i X_j$.

\dfn{Matrice di correlazione}{
  La \term{matrice di correlazione} del vettore $\vett{X}$ è la matrice $n \times n$
  \[ \matr{R}_{\vett{X}} \stackrel{\triangle}{=} E\left[ \vett{X}\vett{X}^T \right] = \begin{bmatrix} R_{X_1 X_1} &
     R_{X_1 X_2} & \cdots & R_{X_1 X_n} \\ R_{X_2 X_1} & R_{X_2 X_2} & \cdots & R_{X_2 X_n} \\ \vdots & \vdots & \ddots &
     \vdots \\ R_{X_n X_1} & R_{X_n X_2} & \cdots & R_{X_n X_n} \end{bmatrix} , \qquad R_{X_i X_j} = E[X_i X_j] . \]
}

È \textbf{simmetrica}, perché $E[X_i X_j] = E[X_j X_i]$, e sulla diagonale principale ci sono i \textbf{valori quadratici
medi} delle variabili, perché $E[X_i X_i] = E[X_i^2]$.

\dfn{Matrice di covarianza}{
  La \term{matrice di covarianza} del vettore $\vett{X}$ è la matrice $n \times n$
  \[ \matr{C}_{\vett{X}} \stackrel{\triangle}{=} E\left[ \left( \vett{X} - \vett{\mu}_{\vett{X}} \right) \left( \vett{X}
     - \vett{\mu}_{\vett{X}} \right)^T \right] = \begin{bmatrix} C_{X_1 X_1} & C_{X_1 X_2} & \cdots & C_{X_1 X_n} \\
     C_{X_2 X_1} & C_{X_2 X_2} & \cdots & C_{X_2 X_n} \\ \vdots & \vdots & \ddots & \vdots \\ C_{X_n X_1} & C_{X_n X_2} &
     \cdots & C_{X_n X_n} \end{bmatrix} , \qquad C_{X_i X_j} = COV(X_i, X_j) . \]
}

Anche questa è simmetrica, e l'elemento $i$-esimo della diagonale principale è
$C_{X_i X_i} = E[(X_i - \mu_{X_i})(X_i - \mu_{X_i})] = \sigma_{X_i}^2$, la \textbf{varianza} di $X_i$. Fuori dalla
diagonale ci sono le covarianze, che si possono scrivere con i coefficienti di correlazione:
$C_{X_i X_j} = \rho_{X_i X_j}\, \sigma_{X_i} \sigma_{X_j}$. Per $n = 2$, scrivendo in breve $\sigma_i = \sigma_{X_i}$ e
$\rho = \rho_{X_1 X_2}$:
\[ \matr{C}_{\vett{X}} = \begin{bmatrix} \sigma_1^2 & \rho\, \sigma_1 \sigma_2 \\ \rho\, \sigma_1 \sigma_2 & \sigma_2^2
   \end{bmatrix} . \]

\teorema[covcorr]{Covarianza e correlazione}{
  \[ \matr{C}_{\vett{X}} = \matr{R}_{\vett{X}} - \vett{\mu}_{\vett{X}} \vett{\mu}_{\vett{X}}^T . \]
}

\begin{dimostrazione}[come nel Capitolo~9, ma tutto insieme]
\passo{Si sviluppa il prodotto} Scriviamo in breve $\vett{\mu}$ per $\vett{\mu}_{\vett{X}}$. Il prodotto fra matrici si
sviluppa come quello fra numeri, facendo attenzione a non scambiare l'ordine dei fattori; inoltre
$(\vett{X} - \vett{\mu})^T = \vett{X}^T - \vett{\mu}^T$:
\[ (\vett{X} - \vett{\mu})(\vett{X} - \vett{\mu})^T = \vett{X}\vett{X}^T - \vett{X}\vett{\mu}^T -
   \vett{\mu}\vett{X}^T + \vett{\mu}\vett{\mu}^T . \]
\passo{Si fa la media} La media è lineare e si fa elemento per elemento; $\vett{\mu}$ è costante ed esce dalla media:
$E[\vett{X}\vett{\mu}^T] = E[\vett{X}]\, \vett{\mu}^T = \vett{\mu}\vett{\mu}^T$, e allo stesso modo
$E[\vett{\mu}\vett{X}^T] = \vett{\mu}\vett{\mu}^T$. Quindi
\[ \matr{C}_{\vett{X}} = \matr{R}_{\vett{X}} - \vett{\mu}\vett{\mu}^T - \vett{\mu}\vett{\mu}^T + \vett{\mu}\vett{\mu}^T
   = \matr{R}_{\vett{X}} - \vett{\mu}\vett{\mu}^T . \]
Elemento per elemento è proprio la formula $COV(X_i, X_j) = E[X_i X_j] - \mu_{X_i} \mu_{X_j}$ del Capitolo~9.
\end{dimostrazione}

\subsection{La caratterizzazione sintetica}

Per una sola variabile la coppia media-varianza $(\mu_X, \sigma_X^2)$ riassume spesso le informazioni più importanti;
per una coppia, la «cinquina» $(\mu_X, \sigma_X^2, \mu_Y, \sigma_Y^2, COV(X, Y))$. Per un vettore aleatorio lo stesso
ruolo è svolto dal vettore delle medie $\vett{\mu}_{\vett{X}}$ insieme alla matrice di covarianza $\matr{C}_{\vett{X}}$
(o a quella di correlazione $\matr{R}_{\vett{X}}$, che per il Teorema~\ref{teo:covcorr} contiene la stessa
informazione, noto $\vett{\mu}_{\vett{X}}$): è la \term{caratterizzazione statistica sintetica} del vettore. Per
$n = 2$ si ritrova proprio la cinquina: due medie e, nella matrice $2 \times 2$ simmetrica, due varianze e una
covarianza. In generale servono $n$ medie e $\frac{n(n+1)}{2}$ elementi della matrice (la diagonale e il triangolo sopra
di essa; quello sotto è uguale per simmetria).

La caratterizzazione sintetica è molto più povera di quella completa, la densità congiunta $f_{\vett{X}}$: dalla densità
si calcolano medie e covarianze, ma in generale non si può tornare indietro, perché vettori con densità molto diverse
possono avere le stesse medie e covarianze. Vedremo che per i vettori gaussiani, invece, si può.

\begin{esempio}{Il vettore della densità $x + y$}
Riprendiamo la coppia del Capitolo~14 (là si chiamavano $X$ e $Y$), con densità $f_{\vett{X}}(\vett{x}) = x_1 + x_2$ sul
quadrato $[0, 1]^2$. Avevamo trovato $E[X_1] = E[X_2] = \frac{7}{12}$, $E[X_1^2] = E[X_2^2] = \frac{5}{12}$ e
$E[X_1 X_2] = \frac13$. Quindi
\[ \vett{\mu}_{\vett{X}} = \begin{bmatrix} 7/12 \\ 7/12 \end{bmatrix} , \qquad
   \matr{R}_{\vett{X}} = \begin{bmatrix} 5/12 & 1/3 \\ 1/3 & 5/12 \end{bmatrix} = \frac{1}{144} \begin{bmatrix} 60 & 48
   \\ 48 & 60 \end{bmatrix} , \qquad
   \vett{\mu}_{\vett{X}} \vett{\mu}_{\vett{X}}^T = \frac{1}{144} \begin{bmatrix} 49 & 49 \\ 49 & 49 \end{bmatrix} , \]
dove $\vett{\mu}_{\vett{X}} \vett{\mu}_{\vett{X}}^T$ è il prodotto colonna per riga, con tutti gli elementi uguali a
$\frac{7}{12} \cdot \frac{7}{12} = \frac{49}{144}$. Per il Teorema~\ref{teo:covcorr}
\[ \matr{C}_{\vett{X}} = \matr{R}_{\vett{X}} - \vett{\mu}_{\vett{X}} \vett{\mu}_{\vett{X}}^T = \frac{1}{144}
   \begin{bmatrix} 11 & -1 \\ -1 & 11 \end{bmatrix} . \]
Sulla diagonale ci sono le varianze $\frac{11}{144}$, fuori la covarianza $-\frac{1}{144}$: gli stessi valori del
Capitolo~14, ottenuti tutti insieme.
\end{esempio}

% ============================================================
\section{La matrice di covarianza è sempre definita non negativa}
% ============================================================

Una varianza non può essere negativa. Per una matrice di covarianza vale qualcosa di più forte, che riguarda tutte le
combinazioni lineari delle variabili del vettore.

\dfn{Matrice definita non negativa}{
  Una matrice quadrata $\matr{C}$ di dimensioni $M \times M$ si dice \term{definita non negativa} se, per qualunque
  vettore $\vett{q} \in \R^M$ (cioè $M \times 1$),
  \[ \vett{q}^T \matr{C}\, \vett{q} \geq 0 . \]
}

\teorema[nonnegativa]{La matrice di covarianza è definita non negativa}{
  La matrice di covarianza $\matr{C}_{\vett{X}}$ di qualunque vettore aleatorio $\vett{X}$ di dimensione $M$ è definita
  non negativa: $\vett{q}^T \matr{C}_{\vett{X}}\, \vett{q} \geq 0$ per ogni $\vett{q} \in \R^M$.
}

\begin{dimostrazione}[la forma quadratica è una varianza]
\passo{Una combinazione lineare} Fissato $\vett{q}$, la quantità $\vett{q}^T \vett{X} = q_1 X_1 + \cdots + q_M X_M$ è una
variabile aleatoria scalare (riga per colonna dà un numero). Per la linearità della media,
$E[\vett{q}^T \vett{X}] = \vett{q}^T \vett{\mu}_{\vett{X}}$.
\passo{Il suo valore quadratico medio} Un numero è uguale al suo trasposto, quindi
$\vett{q}^T \vett{X} = (\vett{q}^T \vett{X})^T = \vett{X}^T \vett{q}$ (trasposto di un prodotto, con l'ordine
invertito). Allora
\[ E[(\vett{q}^T \vett{X})^2] = E[(\vett{q}^T \vett{X})(\vett{X}^T \vett{q})] = E[\vett{q}^T \vett{X}\vett{X}^T
   \vett{q}] = \vett{q}^T E[\vett{X}\vett{X}^T]\, \vett{q} , \]
dove $\vett{q}$, che è costante, esce dalla media.
\passo{La varianza non è negativa}
\[ 0 \leq \mathrm{VAR}(\vett{q}^T \vett{X}) = E[(\vett{q}^T \vett{X})^2] - \left( E[\vett{q}^T \vett{X}] \right)^2 =
   \vett{q}^T E[\vett{X}\vett{X}^T]\, \vett{q} - (\vett{q}^T \vett{\mu}_{\vett{X}})^2 . \]
(Nella slide 77 manca il quadrato su $E[\vett{q}^T \vett{X}]$: la varianza è il valore quadratico medio meno il
\emph{quadrato} della media.)
\passo{Anche il quadrato della media è una forma quadratica} Con lo stesso trucco del passo 2,
$(\vett{q}^T \vett{\mu}_{\vett{X}})^2 = (\vett{q}^T \vett{\mu}_{\vett{X}})(\vett{\mu}_{\vett{X}}^T \vett{q}) =
\vett{q}^T \vett{\mu}_{\vett{X}} \vett{\mu}_{\vett{X}}^T \vett{q}$. Raccogliendo $\vett{q}^T$ a sinistra e $\vett{q}$ a
destra:
\[ 0 \leq \vett{q}^T \left( E[\vett{X}\vett{X}^T] - \vett{\mu}_{\vett{X}} \vett{\mu}_{\vett{X}}^T \right) \vett{q} =
   \vett{q}^T \matr{C}_{\vett{X}}\, \vett{q} , \]
per il Teorema~\ref{teo:covcorr}. E questo vale qualunque sia $\vett{q} \in \R^M$.
\end{dimostrazione}

Il significato è semplice: $\vett{q}^T \matr{C}_{\vett{X}} \vett{q}$ \textbf{è la varianza della combinazione lineare}
$\vett{q}^T \vett{X}$, e una varianza non può essere negativa. Con $M = 2$, svolgendo i prodotti riga per colonna:
\[ \vett{q}^T \matr{C}_{\vett{X}} \vett{q} = [q_1, q_2] \begin{bmatrix} \sigma_1^2 & \rho\sigma_1\sigma_2 \\
   \rho\sigma_1\sigma_2 & \sigma_2^2 \end{bmatrix} \begin{bmatrix} q_1 \\ q_2 \end{bmatrix} = [q_1, q_2]
   \begin{bmatrix} \sigma_1^2 q_1 + \rho\sigma_1\sigma_2 q_2 \\ \rho\sigma_1\sigma_2 q_1 + \sigma_2^2 q_2 \end{bmatrix}
   = q_1^2 \sigma_1^2 + q_2^2 \sigma_2^2 + 2 q_1 q_2\, \rho\, \sigma_1 \sigma_2 , \]
che è la formula della varianza di $q_1 X_1 + q_2 X_2$ del Capitolo~9 (Teorema~\ref{teo:varsomma}).

\info{Un vecchio risultato in veste nuova}{
  Scegliendo $\vett{q} = [1/\sigma_1, \pm 1/\sigma_2]^T$ la forma quadratica vale $1 + 1 \pm 2\rho$, e la condizione
  $\geq 0$ diventa $\rho \geq -1$ (con il segno $+$) e $\rho \leq 1$ (con il segno $-$): è la dimostrazione di
  $\abs{\rho} \leq 1$ del Capitolo~9. La non negatività della matrice di covarianza generalizza quel risultato a $M$
  variabili: nessuna combinazione delle variabili può avere varianza negativa.
}

% ============================================================
\section{Vettori gaussiani}
% ============================================================

\subsection{Si parte da \texorpdfstring{$n$}{n} gaussiane indipendenti}

Il punto di partenza è il caso più semplice: $n$ variabili gaussiane statisticamente indipendenti,
$X_i \sim \mathcal{N}(\mu_{X_i}, \sigma_{X_i}^2)$. La loro densità congiunta si può scrivere in una forma che usa solo il
vettore delle medie e la matrice di covarianza.

\teorema[gaussindip]{Densità di $n$ gaussiane indipendenti}{
  Se $X_1, \ldots, X_n$ sono gaussiane indipendenti, la densità congiunta del vettore $\vett{X} = [X_1, \ldots, X_n]^T$ è
  \[ f_{\vett{X}}(\vett{x}) = \frac{1}{\sqrt{(2\pi)^n \det \matr{C}_{\vett{X}}}}\; e^{-\frac12 \left( \vett{x} -
     \vett{\mu}_{\vett{X}} \right)^T \matr{C}_{\vett{X}}^{-1} \left( \vett{x} - \vett{\mu}_{\vett{X}} \right)} . \]
}

\begin{dimostrazione}
\passo{Il prodotto delle marginali} Per l'indipendenza la densità congiunta è il prodotto delle densità:
\[ f_{\vett{X}}(\vett{x}) = \prod_{i=1}^{n} f_{X_i}(x_i) = \prod_{i=1}^{n} \frac{1}{\sqrt{2\pi\sigma_{X_i}^2}}\,
   e^{-\frac{(x_i - \mu_{X_i})^2}{2\sigma_{X_i}^2}} = \frac{1}{\sqrt{(2\pi)^n \prod_{i=1}^{n} \sigma_{X_i}^2}}\;
   e^{-\frac12 \sum_{i=1}^{n} \frac{(x_i - \mu_{X_i})^2}{\sigma_{X_i}^2}} . \]
Si sono usate due regole: il prodotto delle radici è la radice del prodotto ($\sqrt{a}\sqrt{b} = \sqrt{ab}$), e il
prodotto di esponenziali è l'esponenziale della somma ($e^a e^b = e^{a + b}$).
\passo{La matrice di covarianza è diagonale} Le variabili sono indipendenti, quindi incorrelate a due a due: tutte le
covarianze fuori dalla diagonale sono nulle. Per una matrice diagonale determinante e inversa sono immediati:
\[ \matr{C}_{\vett{X}} = \begin{bmatrix} \sigma_{X_1}^2 & 0 & \cdots & 0 \\ 0 & \sigma_{X_2}^2 & \cdots & 0 \\ \vdots
   & \vdots & \ddots & \vdots \\ 0 & 0 & \cdots & \sigma_{X_n}^2 \end{bmatrix} \quad\Longrightarrow\quad
   \det \matr{C}_{\vett{X}} = \prod_{i=1}^{n} \sigma_{X_i}^2 , \qquad
   \matr{C}_{\vett{X}}^{-1} = \begin{bmatrix} 1/\sigma_{X_1}^2 & \cdots & 0 \\ \vdots & \ddots & \vdots \\ 0 & \cdots &
   1/\sigma_{X_n}^2 \end{bmatrix} . \]
Il prodotto delle varianze sotto la radice è quindi $\det \matr{C}_{\vett{X}}$.
\passo{L'esponente è una forma quadratica} La somma all'esponente è una somma di quadrati pesati con i pesi
$\frac{1}{\sigma_{X_i}^2}$, cioè la forma quadratica della matrice diagonale $\matr{C}_{\vett{X}}^{-1}$ calcolata nel
vettore $\vett{x} - \vett{\mu}_{\vett{X}}$:
\[ \sum_{i=1}^{n} \frac{(x_i - \mu_{X_i})^2}{\sigma_{X_i}^2} = \left( \vett{x} - \vett{\mu}_{\vett{X}} \right)^T
   \matr{C}_{\vett{X}}^{-1} \left( \vett{x} - \vett{\mu}_{\vett{X}} \right) . \]
\end{dimostrazione}

\subsection{La definizione generale}

La formula appena trovata non contiene più le singole variabili: solo $\vett{\mu}_{\vett{X}}$ e $\matr{C}_{\vett{X}}$.
L'idea è usarla anche quando la matrice di covarianza \textbf{non} è diagonale.

\dfn{Vettore gaussiano}{
  Un vettore aleatorio $\vett{X}$ si dice \term{gaussiano}, e si scrive
  $\vett{X} \sim \mathcal{N}(\vett{\mu}_{\vett{X}}, \matr{C}_{\vett{X}})$, quando la sua densità congiunta ha la forma
  \[ f_{\vett{X}}(\vett{x}) = \frac{1}{\sqrt{(2\pi)^n \det \matr{C}_{\vett{X}}}}\; e^{-\frac12 \left( \vett{x} -
     \vett{\mu}_{\vett{X}} \right)^T \matr{C}_{\vett{X}}^{-1} \left( \vett{x} - \vett{\mu}_{\vett{X}} \right)} , \]
  \textbf{anche se le componenti non sono indipendenti}, cioè anche se $\matr{C}_{\vett{X}}$ non è diagonale.
  Equivalentemente, $n$ variabili aleatorie con una densità congiunta di questa forma si dicono
  \term{congiuntamente gaussiane}.
}

Perché la formula abbia senso, $\matr{C}_{\vett{X}}$ deve essere invertibile ($\det \matr{C}_{\vett{X}} \neq 0$; per il
Teorema~\ref{teo:nonnegativa}, anzi, $\det \matr{C}_{\vett{X}} > 0$). È la naturale generalizzazione della densità di una
gaussiana: con $n = 1$ la matrice di covarianza è il numero $\sigma_X^2$, il suo determinante è $\sigma_X^2$ e la sua
inversa è $\frac{1}{\sigma_X^2}$, e la formula ridà $\frac{1}{\sqrt{2\pi\sigma_X^2}}\, e^{-\frac{(x -
\mu_X)^2}{2\sigma_X^2}}$.
\begin{center}
\renewcommand{\arraystretch}{1.6}
\begin{tabular}{@{}lcc@{}}
  \toprule
  & una variabile, $\mathcal{N}(\mu_X, \sigma_X^2)$ & vettore di $n$ variabili, $\mathcal{N}(\vett{\mu}_{\vett{X}},
  \matr{C}_{\vett{X}})$ \\ \midrule
  centro & $\mu_X$ & $\vett{\mu}_{\vett{X}}$ \\
  larghezza & $\sigma_X^2$ & $\matr{C}_{\vett{X}}$ \\
  costante davanti & $\dfrac{1}{\sqrt{2\pi\, \sigma_X^2}}$ & $\dfrac{1}{\sqrt{(2\pi)^n \det \matr{C}_{\vett{X}}}}$ \\
  esponente & $-\frac12\, (x - \mu_X)\, \frac{1}{\sigma_X^2}\, (x - \mu_X)$ & $-\frac12 (\vett{x} -
  \vett{\mu}_{\vett{X}})^T\, \matr{C}_{\vett{X}}^{-1} (\vett{x} - \vett{\mu}_{\vett{X}})$ \\ \bottomrule
\end{tabular}
\end{center}
La forma quadratica all'esponente misura quanto $\vett{x}$ è lontano dal centro $\vett{\mu}_{\vett{X}}$, tenendo conto
delle larghezze e delle correlazioni: è il quadrato di una «distanza», e la densità è tanto più alta quanto più
$\vett{x}$ è vicino al centro. I punti alla stessa distanza, cioè con la stessa densità, formano delle \textbf{ellissi}
(in tre dimensioni, degli ellissoidi).

\subsection{Il caso di due variabili}

Con $n = 2$, $\vett{\mu}_{\vett{X}} = [\mu_1, \mu_2]^T$ e la matrice di covarianza scritta con il coefficiente di
correlazione, la formula dell'inversa $2 \times 2$ del richiamo dà
\[ \matr{C}_{\vett{X}} = \begin{bmatrix} \sigma_1^2 & \rho\sigma_1\sigma_2 \\ \rho\sigma_1\sigma_2 & \sigma_2^2
   \end{bmatrix} , \quad \det \matr{C}_{\vett{X}} = \sigma_1^2 \sigma_2^2 (1 - \rho^2) , \quad
   \matr{C}_{\vett{X}}^{-1} = \frac{1}{\sigma_1^2 \sigma_2^2 (1 - \rho^2)} \begin{bmatrix} \sigma_2^2 &
   -\rho\sigma_1\sigma_2 \\ -\rho\sigma_1\sigma_2 & \sigma_1^2 \end{bmatrix} . \]
Svolgendo i prodotti all'esponente (come per la forma quadratica della sezione precedente) si ottiene la
\term{gaussiana bivariata}:
\[ f_{X_1 X_2}(x_1, x_2) = \frac{1}{2\pi \sigma_1 \sigma_2 \sqrt{1 - \rho^2}}\, \exp\left\{ -\frac{1}{2(1 - \rho^2)}
   \left[ \frac{(x_1 - \mu_1)^2}{\sigma_1^2} - 2\rho\, \frac{(x_1 - \mu_1)(x_2 - \mu_2)}{\sigma_1 \sigma_2} +
   \frac{(x_2 - \mu_2)^2}{\sigma_2^2} \right] \right\} , \]
dove $\exp\{a\}$ è un altro modo di scrivere $e^a$, comodo quando l'esponente è lungo. Serve $\abs{\rho} < 1$: con
$\rho = \pm 1$ il determinante si annulla, la matrice non è invertibile e le due variabili sono legate da una relazione
lineare esatta (tutta la probabilità sta su una retta, e una densità nel piano non esiste).

\begin{figure}[H]
\centering
\begin{tikzpicture}
\begin{groupplot}[group style={group size=2 by 1, horizontal sep=1.6cm}, width=8cm, height=6.6cm, view={-35}{38},
  xmin=-3, xmax=3, ymin=-3, ymax=3, zmin=0, zmax=0.27, xtick={-2,0,2}, ytick={-2,0,2}, ztick=\empty,
  xlabel=$x_1$, ylabel=$x_2$, tick label style={font=\scriptsize}, label style={font=\small}, colormap/viridis]
\nextgroupplot[title={$\rho = 0$}]
  \addplot3[surf, shader=faceted, opacity=0.6, draw=black!30, line width=0.2pt, domain=-3:3, y domain=-3:3,
    samples=31] {exp(-(x^2 + y^2)/2)/(2*pi)};
\nextgroupplot[title={$\rho = 0{,}8$}]
  \addplot3[surf, shader=faceted, opacity=0.6, draw=black!30, line width=0.2pt, domain=-3:3, y domain=-3:3,
    samples=31] {exp(-(x^2 - 1.6*x*y + y^2)/(2*0.36))/(2*pi*0.6)};
\end{groupplot}
\end{tikzpicture}
\caption{La gaussiana bivariata con $\mu_1 = \mu_2 = 0$ e $\sigma_1 = \sigma_2 = 1$. Con $\rho = 0$ la campana è
rotonda; con $\rho = 0{,}8$ è allungata lungo la bisettrice $x_2 = x_1$, e più alta, perché il volume sotto la superficie
resta 1.}
\end{figure}

\begin{figure}[H]
\centering
\begin{tikzpicture}
\begin{groupplot}[group style={group size=3 by 1, horizontal sep=0.45cm, yticklabels at=edge left}, width=6.6cm,
  height=6.6cm, xmin=-3.3, xmax=3.3, ymin=-3.3, ymax=3.3, axis equal image, xtick={-2,0,2}, ytick={-2,0,2},
  tick label style={font=\scriptsize}, label style={font=\small}, xlabel=$x_1$, grid=major, grid style={mspigray!20}]
\nextgroupplot[title={$\rho = 0$}, ylabel=$x_2$]
  \addplot[only marks, mark=*, mark size=0.7pt, mspiA!75] coordinates {(-0.74,0.87) (-0.94,0.54) (-1.23,0.79) (-0.08,1.42) (0.85,0.05) (-0.28,0.48) (-0.04,-0.86) (0.61,-2.43) (-0.05,2.12) (0.53,-1.39) (1.09,0.19) (-0.59,-0.60) (0.14,0.85) (-0.09,1.76) (-0.60,0.86) (-1.52,-1.91) (-0.08,0.29) (1.87,-0.31) (0.31,-1.18) (0.06,0.58) (-0.37,1.05) (1.42,0.62) (-0.02,-0.16) (0.64,-0.97) (0.99,-0.55) (0.31,-1.06) (-0.89,-0.77) (0.55,0.08) (0.11,1.01) (-0.85,0.91) (-0.54,-0.85) (-0.51,1.08) (1.67,-0.51) (-1.03,-1.53) (0.06,0.46) (0.42,0.28) (1.40,0.08) (0.43,-0.24) (-0.92,-2.19) (0.06,-0.57) (-0.50,-0.59) (0.07,0.13) (-2.21,-0.36) (0.45,0.34) (-0.21,0.18) (-1.58,0.36) (-0.27,0.07) (0.65,0.13) (-1.13,-1.97) (-0.47,-0.01) (0.22,1.01) (0.39,0.31) (-1.75,0.77) (0.16,2.06) (0.76,-0.75) (0.10,0.63) (-0.93,-1.62) (0.44,-0.53) (-0.19,0.09) (-0.28,0.79) (0.28,0.28) (0.85,-0.04) (-1.24,0.01) (-1.31,0.23) (-1.16,-0.02) (1.24,0.68) (-0.78,-1.40) (-0.43,1.03) (0.56,-1.51) (0.85,0.29) (0.75,1.39) (-0.81,0.57) (0.19,1.08) (-0.28,-1.38) (0.18,0.99) (-0.24,-0.24) (0.18,-0.12) (-0.40,-0.40) (1.48,-0.42) (0.20,1.17) (1.11,-1.62) (-0.74,-1.22) (-0.11,-2.22) (-1.38,-0.12) (0.36,-1.14) (0.26,1.29) (-0.92,1.28) (-0.03,-2.18) (0.17,0.13) (1.47,-0.68) (1.83,-0.03) (0.56,-0.17) (0.54,1.27) (0.02,-0.50) (0.68,-0.23) (0.31,1.03) (0.18,0.39) (0.23,0.12) (0.88,1.22) (0.32,-1.69) (-1.45,0.83) (1.50,1.32) (-0.43,0.27) (-0.49,0.44) (-0.04,-0.97) (0.65,0.07) (2.16,1.19) (-1.48,0.93) (-0.86,-0.57) (-0.14,-1.26) (3.08,0.86) (-0.85,1.18) (-2.22,1.28) (0.17,0.30) (0.98,0.27) (0.63,-0.89) (0.50,-0.23) (0.59,-0.13) (0.49,-0.44) (-1.20,-0.68) (1.01,-0.19) (0.18,-0.27) (0.97,-0.72) (-2.56,-0.61) (-0.54,1.27) (-0.73,-0.44) (0.64,0.84) (-0.28,1.34) (-0.35,0.23) (-0.84,-2.11) (-0.62,-0.56) (-0.82,2.38) (-1.51,0.55) (1.64,0.83) (0.61,-0.48) (-0.51,1.36) (1.68,-0.51) (0.49,0.07) (0.78,-0.65) (1.05,-0.32) (0.36,-0.60) (1.11,-0.25) (0.66,0.02) (0.91,0.52) (2.57,-0.71) (0.92,-0.45) (0.20,-1.64) (1.29,-2.30) (0.85,-1.48) (-0.25,0.56) (-0.72,0.41) (0.01,-0.58) (0.47,1.55) (0.41,0.10) (0.20,0.00) (1.34,0.95) (-0.06,0.47) (-0.94,0.01) (0.37,0.01) (-0.03,1.04) (0.77,0.70) (-2.28,0.23) (-0.66,-0.54) (0.11,0.63) (-1.88,1.62) (0.56,-0.11) (-1.04,-1.34) (1.50,-0.76) (-0.29,-0.32) (-0.21,2.36) (1.76,-1.11) (-0.91,1.49) (0.24,-0.98) (-0.57,0.95) (-0.19,-0.05) (-0.86,0.57) (-1.04,-0.65) (-0.40,-1.17) (0.30,0.04) (-0.08,0.94) (1.17,-0.10) (1.29,1.39) (-0.41,0.96) (-1.07,-1.18) (0.04,1.03) (2.33,0.14) (1.32,-1.30) (0.39,1.35) (-0.52,0.24) (0.47,1.97) (-1.05,-0.85) (0.87,1.00) (0.53,-0.42) (0.22,2.32) (-0.17,1.10) (1.42,-0.81) (-0.35,0.83) (0.01,0.29) (-0.56,-0.08) (-0.08,1.87) (0.30,0.28) (1.34,0.45) (0.02,0.79) (-1.61,-0.21) (0.34,-0.07) (-0.10,0.14) (-0.89,1.18) (-0.56,0.90) (-0.97,2.63) (0.77,-1.95) (-1.60,-1.39) (-0.27,-0.35) (1.66,1.81) (-1.55,-0.53) (-1.98,-0.77) (0.18,0.56) (0.32,1.59) (-1.01,-0.70) (0.30,1.01) (-0.11,0.44)};
  \pgfplotsinvokeforeach{1,2,3}{
    \addplot[mspiB, line width=0.9pt, domain=0:360, samples=73]
      ({#1*(0.7071*cos(x) + 0.7071*sin(x))}, {#1*(0.7071*cos(x) - 0.7071*sin(x))});
  }
\nextgroupplot[title={$\rho = 0{,}8$}]
  \addplot[only marks, mark=*, mark size=0.7pt, mspiA!75] coordinates {(1.72,1.07) (0.26,-0.35) (0.96,-0.16) (0.26,1.69) (0.12,-0.09) (0.34,0.43) (-0.28,-0.53) (0.16,0.25) (-0.79,-0.43) (-0.58,-0.61) (-0.44,-1.15) (-2.05,-0.69) (-1.05,-1.25) (0.07,0.56) (-0.19,-0.69) (0.81,1.04) (-0.55,-0.42) (0.09,1.44) (0.40,-0.28) (0.08,-0.80) (0.42,1.05) (1.25,1.33) (-0.45,-0.43) (-0.74,-0.64) (-1.61,-2.80) (0.13,-0.11) (-1.04,-1.26) (0.71,0.32) (0.79,0.96) (-1.03,0.15) (1.01,0.28) (0.07,-0.58) (-0.15,0.51) (-0.37,-0.78) (0.52,-0.17) (0.06,0.76) (0.64,0.09) (-0.48,0.27) (1.70,1.03) (-0.14,-0.56) (-1.14,-1.49) (-0.37,0.40) (1.20,2.08) (-0.71,-1.32) (0.70,-0.37) (-0.76,0.08) (1.07,1.50) (-0.85,-0.15) (0.40,-0.12) (0.67,1.53) (-0.54,0.82) (0.53,0.28) (-2.22,-2.31) (-0.83,-1.36) (0.69,1.04) (0.01,-0.01) (-0.21,-0.67) (-0.41,0.79) (-0.80,-0.12) (0.85,0.67) (0.85,0.55) (-0.28,-0.33) (1.11,0.79) (-2.56,-2.00) (-0.19,-0.09) (-0.11,0.02) (1.97,1.87) (0.35,-0.65) (-0.62,-0.95) (-0.76,-0.83) (-1.21,-1.72) (-0.58,-0.14) (-0.35,0.09) (1.67,0.80) (-0.17,0.13) (-1.68,-1.35) (-0.42,0.65) (-0.47,-0.09) (0.57,-0.43) (-1.19,-0.35) (-0.63,0.44) (-0.72,-1.22) (-0.48,0.41) (0.12,1.69) (-1.56,-1.62) (0.47,0.48) (-0.33,0.08) (-0.99,-0.23) (0.07,-0.24) (-0.26,-0.47) (0.07,-1.13) (-0.27,-0.75) (-0.93,0.04) (0.30,0.00) (-0.17,1.08) (0.12,-0.32) (0.18,0.30) (-0.42,0.80) (-0.56,-0.60) (0.26,0.44) (-0.41,-0.12) (1.00,1.27) (0.08,0.14) (1.69,1.60) (-0.07,-0.12) (1.90,1.19) (2.08,1.74) (0.70,0.56) (0.90,1.25) (-0.38,-0.66) (-1.49,-1.75) (0.56,0.59) (0.02,-0.04) (-0.23,0.27) (-0.14,-0.16) (1.36,1.10) (-0.06,-0.25) (-0.52,-0.69) (0.57,-0.35) (-0.52,-0.73) (-2.68,-1.73) (-1.48,-1.24) (0.07,-0.43) (-0.94,0.77) (-0.18,-0.50) (-0.86,-0.79) (-0.71,-1.36) (0.79,0.02) (0.81,1.19) (0.87,1.12) (-0.83,-0.38) (0.01,0.37) (-0.16,0.75) (0.56,1.05) (2.51,1.60) (0.03,-0.70) (-1.00,-1.14) (0.80,0.20) (1.97,1.99) (0.39,1.61) (0.63,0.19) (-0.05,1.16) (0.91,0.29) (0.95,1.18) (0.86,0.44) (-1.28,-2.04) (-0.89,-0.85) (-0.69,-0.81) (0.33,-0.10) (0.12,-0.63) (0.33,1.55) (0.53,0.31) (-1.12,0.18) (1.09,1.60) (-0.34,0.37) (-0.11,0.09) (-2.23,-1.76) (0.41,1.23) (1.61,2.11) (0.73,0.64) (0.03,0.87) (-1.48,-1.28) (-1.79,-1.95) (0.25,0.59) (-1.60,-1.68) (-2.02,-1.13) (0.02,1.26) (-1.61,-1.33) (2.11,1.52) (0.43,-0.53) (-0.94,0.03) (-1.01,-1.30) (1.47,2.06) (0.86,-0.65) (-0.90,-0.62) (-1.83,-2.14) (-0.67,-0.13) (-1.96,-3.28) (-1.30,-1.78) (-0.37,-0.22) (0.88,0.26) (-0.45,-0.78) (0.25,-0.11) (-2.12,-2.05) (-0.02,-0.23) (1.61,0.06) (0.69,0.39) (1.26,1.02) (-0.32,-0.52) (-0.33,-0.25) (1.83,2.23) (-0.61,-0.81) (-0.49,-0.55) (-0.98,-0.06) (-0.16,0.12) (0.43,0.26) (-0.80,-1.04) (-1.59,-1.50) (-0.10,0.57) (-0.49,0.18) (-0.72,-1.03) (0.31,-0.63) (-0.40,-0.87) (-0.35,0.47) (-0.75,-1.13) (0.65,0.50) (-2.67,-2.45) (-0.23,0.79) (-0.56,-1.72) (-0.64,-0.18) (-0.14,-1.11) (-0.92,-1.21) (0.98,1.13) (-0.95,0.06) (1.17,0.85) (0.44,-0.17) (-1.28,-0.59) (0.45,0.98) (0.88,0.89) (2.12,0.48)};
  \pgfplotsinvokeforeach{1,2,3}{
    \addplot[mspiB, line width=0.9pt, domain=0:360, samples=73]
      ({#1*(0.9487*cos(x) + 0.3162*sin(x))}, {#1*(0.9487*cos(x) - 0.3162*sin(x))});
  }
\nextgroupplot[title={$\rho = -0{,}6$}]
  \addplot[only marks, mark=*, mark size=0.7pt, mspiA!75] coordinates {(0.02,1.03) (-0.24,0.12) (0.81,0.50) (1.45,0.29) (0.69,-0.73) (0.39,-0.19) (1.82,-0.24) (-1.93,-0.54) (-1.06,-0.47) (-1.34,1.24) (0.53,-0.49) (0.23,-0.04) (-0.61,0.83) (-0.32,0.98) (-1.11,0.53) (-0.71,1.96) (0.73,0.82) (-0.07,-0.15) (-1.72,2.16) (-0.42,0.24) (0.12,-0.54) (-1.63,0.73) (2.88,-1.68) (0.16,-0.03) (-0.02,0.20) (-0.62,-0.32) (-1.31,0.70) (0.79,-0.07) (1.97,-1.70) (-1.11,0.37) (0.93,-2.16) (1.02,-0.44) (1.67,-1.01) (-0.87,-0.48) (1.47,-0.44) (-0.36,0.17) (0.45,0.13) (-0.43,0.71) (-0.93,-0.45) (0.50,-1.30) (-0.97,0.03) (-0.55,-0.87) (1.60,-0.30) (2.91,-1.83) (-0.21,0.67) (1.72,-1.57) (1.04,-0.12) (-0.30,-0.03) (0.67,-0.53) (0.47,-0.66) (0.69,-0.32) (1.47,-0.43) (-0.64,0.10) (0.23,-0.74) (-0.52,0.74) (-2.22,0.81) (0.50,-1.76) (0.96,-1.51) (-0.01,1.00) (-1.11,1.12) (-0.41,0.63) (1.00,-1.48) (-1.03,0.17) (0.80,0.20) (1.03,-0.41) (-0.73,1.18) (-0.39,-0.12) (0.41,0.03) (-0.27,1.33) (0.58,0.37) (1.86,-0.72) (1.58,-2.04) (0.23,-0.45) (-0.77,0.68) (-0.86,1.28) (1.39,-1.83) (0.98,0.30) (1.17,-1.20) (0.13,1.22) (0.59,-0.14) (0.91,-0.47) (-1.49,1.30) (0.02,0.28) (-0.43,0.14) (-0.02,-1.02) (-1.00,1.18) (1.22,-0.32) (-2.58,1.31) (2.11,-1.61) (0.69,-0.26) (2.20,-1.46) (1.46,-1.14) (0.48,0.37) (-1.32,1.82) (0.51,0.67) (-2.19,0.46) (0.40,-1.57) (0.05,-0.91) (-0.88,1.18) (0.63,1.43) (1.25,-0.57) (2.23,-2.14) (0.04,0.63) (0.37,2.32) (0.33,-0.95) (1.27,-0.88) (-0.36,1.33) (-1.68,0.09) (-0.59,2.49) (0.09,-0.64) (-0.02,0.37) (0.76,-1.03) (0.95,-1.37) (0.50,-0.20) (-0.28,-0.26) (1.47,-0.16) (1.00,-0.62) (3.22,-2.10) (-1.28,0.76) (-0.19,2.16) (1.43,-1.46) (-0.52,0.86) (2.31,-1.17) (-0.72,0.00) (1.26,-0.20) (1.83,-1.26) (1.14,0.14) (2.58,-1.25) (-0.14,-0.77) (-0.86,-0.07) (-1.33,0.74) (1.19,-1.41) (1.14,-0.80) (0.61,-1.42) (-0.06,-1.37) (-1.96,1.53) (-0.20,0.29) (2.53,-1.52) (-0.58,0.90) (0.41,-0.28) (-0.11,1.43) (-1.33,0.97) (1.63,-0.52) (0.66,0.36) (-2.88,1.77) (-0.37,0.76) (-0.92,1.04) (-0.53,-0.75) (1.01,-0.82) (0.91,-0.93) (0.46,0.72) (-0.55,0.71) (-0.87,-1.31) (-1.20,2.29) (-0.18,-0.52) (1.07,-0.56) (1.05,-1.23) (-1.67,1.48) (-1.31,-0.28) (-1.43,0.70) (-1.40,2.26) (0.45,-0.66) (-0.12,0.14) (2.13,-1.10) (0.86,-1.92) (-0.36,0.81) (1.20,-0.41) (-0.50,1.25) (-1.34,0.62) (-0.76,0.27) (-0.58,0.62) (-0.11,-0.49) (0.57,1.62) (2.34,-2.05) (-1.09,1.05) (-1.13,-0.17) (0.36,-1.38) (0.73,0.93) (-0.67,0.80) (-0.08,0.25) (0.36,-1.33) (-0.01,-1.15) (0.33,-0.02) (-0.42,-1.66) (-1.04,-0.33) (0.75,-0.55) (0.28,0.01) (0.47,-0.23) (0.42,1.40) (0.87,-1.24) (1.01,-0.26) (2.90,-3.10) (0.37,-0.14) (-0.50,1.18) (1.68,-0.77) (-0.59,1.03) (0.62,1.01) (-0.51,0.75) (-1.20,2.09) (1.02,-1.18) (0.67,-1.71) (-0.72,1.79) (-0.12,-0.81) (-0.66,1.24) (-1.23,0.21) (0.43,1.30) (0.39,0.01) (0.22,-0.15) (-1.05,0.16) (-1.12,-0.43) (-0.51,-1.25) (0.78,-0.16) (-0.62,-0.42) (1.33,-0.13) (0.05,-0.17) (0.61,-0.20) (1.39,-1.39) (1.77,-0.66) (0.72,-1.94) (-1.77,1.45)};
  \pgfplotsinvokeforeach{1,2,3}{
    \addplot[mspiB, line width=0.9pt, domain=0:360, samples=73]
      ({#1*(0.4472*cos(x) + 0.8944*sin(x))}, {#1*(0.4472*cos(x) - 0.8944*sin(x))});
  }
\end{groupplot}
\end{tikzpicture}
\caption{220 punti estratti da tre gaussiane bivariate con medie nulle e varianze unitarie (simulazione), e le ellissi
di uguale densità su cui la forma quadratica vale $1$, $4$ e $9$. Con $\rho = 0$ sono circonferenze; con $\rho > 0$ si
allungano lungo $x_2 = x_1$, con $\rho < 0$ lungo $x_2 = -x_1$.}
\end{figure}

Qualunque sia $\rho$, le tre ellissi racchiudono una probabilità del 39\%, dell'86\% e del 99\%: con soli 220 punti le
percentuali dei pannelli si avvicinano a questi valori, senza coincidere. Per $\rho = 0$ lo si verifica con le coordinate
polari, come nell'integrale di Gauss del Capitolo~15: la probabilità di cadere nel cerchio di raggio $k$ è
$\int_0^{2\pi} \int_0^k \frac{1}{2\pi} e^{-r^2/2}\, r \dd r \dd\theta = 1 - e^{-k^2/2}$, che per $k = 1, 2, 3$ vale
$0{,}39$, $0{,}86$ e $0{,}99$.

\subsection{Incorrelazione e indipendenza}

Nel Capitolo~9 abbiamo visto che due variabili indipendenti sono sempre incorrelate, ma che il viceversa in generale è
falso. Per le variabili congiuntamente gaussiane, invece, le due cose coincidono.

\teorema[gaussincorr]{Per le congiuntamente gaussiane, incorrelate vuol dire indipendenti}{
  Se $X_1, \ldots, X_n$ sono congiuntamente gaussiane e incorrelate a due a due, allora sono statisticamente
  indipendenti.
}

\begin{dimostrazione}
\passo{La matrice è diagonale} Se le variabili sono incorrelate, tutte le covarianze fuori dalla diagonale sono nulle e
$\matr{C}_{\vett{X}}$ è diagonale.
\passo{Si legge al contrario il Teorema~\ref{teo:gaussindip}} Con una matrice diagonale la forma quadratica è una somma
di quadrati e il determinante è un prodotto di varianze; ripercorrendo all'indietro i passaggi di quella dimostrazione,
la densità si spezza in un prodotto:
\[ f_{\vett{X}}(\vett{x}) = \prod_{i=1}^{n} \frac{1}{\sqrt{2\pi\sigma_{X_i}^2}}\, e^{-\frac{(x_i - \mu_{X_i})^2}{2
   \sigma_{X_i}^2}} . \]
\passo{I fattori sono le marginali} Integrando rispetto a tutte le variabili tranne $x_i$, ogni altro fattore (una
densità gaussiana, con area 1) dà 1: resta $f_{X_i}(x_i)$, che è quindi proprio l'$i$-esimo fattore. La densità
congiunta è il prodotto delle marginali, cioè le variabili sono indipendenti.
\end{dimostrazione}

\warning{Gaussiane una per una non basta}{
  Il teorema richiede che le variabili siano \emph{congiuntamente} gaussiane, non solo gaussiane una per una. Sia
  $X_1 \sim \mathcal{N}(0, 1)$ e sia $X_2 = S\, X_1$, dove $S$ vale $+1$ o $-1$ con probabilità $\frac12$, indipendente da
  $X_1$. Anche $X_2$ è $\mathcal{N}(0, 1)$: cambiare segno a caso a una campana simmetrica la lascia com'è. Le due
  variabili sono incorrelate, perché $E[X_1 X_2] = E[S]\, E[X_1^2] = 0 \cdot 1 = 0$; eppure sono fortemente dipendenti,
  perché $\abs{X_2} = \abs{X_1}$ sempre. Il vettore $[X_1, X_2]^T$ non è gaussiano: i suoi punti stanno tutti sulle due
  bisettrici $x_2 = \pm x_1$, mentre la densità di un vettore gaussiano con $\abs{\rho} < 1$ è positiva in tutto il piano.
}

\subsection{Marginali e condizionate: si completa il quadrato}

Le slide affermano che le marginali e le condizionate di un vettore gaussiano sono ancora gaussiane. Per due variabili
lo si può verificare con un calcolo elementare. Prendiamo il caso standardizzato, $\mu_1 = \mu_2 = 0$ e
$\sigma_1 = \sigma_2 = 1$ (quello generale si riconduce a questo standardizzando le due variabili), in cui
\[ f_{X_1 X_2}(x_1, x_2) = \frac{1}{2\pi\sqrt{1 - \rho^2}}\, e^{-\frac{x_1^2 - 2\rho x_1 x_2 + x_2^2}{2(1 - \rho^2)}} . \]

\teorema[gaussmarginale]{Marginale e condizionata della gaussiana bivariata}{
  Nel caso standardizzato:
  \begin{enumerate}[label=\alph*)]
    \item la marginale è una gaussiana standard: $X_1 \sim \mathcal{N}(0, 1)$, e lo stesso vale per $X_2$;
    \item la condizionata è gaussiana: sapendo che $X_1 = x_1$, si ha $X_2 \sim \mathcal{N}(\rho x_1,\ 1 - \rho^2)$, cioè
          \[ f_{X_2 \mid X_1}(x_2 \mid x_1) = \frac{1}{\sqrt{2\pi(1 - \rho^2)}}\, e^{-\frac{(x_2 - \rho x_1)^2}{2(1 -
             \rho^2)}} . \]
  \end{enumerate}
}

\begin{dimostrazione}[completando il quadrato]
\passo{Il quadrato} Come quando si risolve un'equazione di secondo grado, si raccolgono i termini che contengono $x_2$
in un quadrato perfetto. Siccome $(x_2 - \rho x_1)^2 = x_2^2 - 2\rho x_1 x_2 + \rho^2 x_1^2$, basta aggiungere e togliere
$\rho^2 x_1^2$:
\[ x_1^2 - 2\rho x_1 x_2 + x_2^2 = (x_2 - \rho x_1)^2 + x_1^2 - \rho^2 x_1^2 = (x_2 - \rho x_1)^2 + (1 - \rho^2)\,
   x_1^2 . \]
\passo{La densità si spezza in due fattori} Dividendo per $2(1 - \rho^2)$, l'esponente diventa
$-\frac{(x_2 - \rho x_1)^2}{2(1 - \rho^2)} - \frac{x_1^2}{2}$; inoltre $2\pi\sqrt{1 - \rho^2} = \sqrt{2\pi} \cdot
\sqrt{2\pi(1 - \rho^2)}$. Quindi
\[ f_{X_1 X_2}(x_1, x_2) = \underbrace{\frac{1}{\sqrt{2\pi}}\, e^{-\frac{x_1^2}{2}}}_{\mathcal{N}(0,\, 1) \text{ in }
   x_1} \cdot \underbrace{\frac{1}{\sqrt{2\pi(1 - \rho^2)}}\, e^{-\frac{(x_2 - \rho x_1)^2}{2(1 -
   \rho^2)}}}_{\mathcal{N}(\rho x_1,\, 1 - \rho^2) \text{ in } x_2} . \]
\passo{La marginale} Per marginalizzare si integra rispetto a $x_2$. Il primo fattore non dipende da $x_2$ ed esce
dall'integrale; il secondo, come funzione di $x_2$, è una densità gaussiana (di media $\rho x_1$), quindi ha area 1:
\[ f_{X_1}(x_1) = \int_{-\infty}^{\infty} f_{X_1 X_2}(x_1, x_2) \dd x_2 = \frac{1}{\sqrt{2\pi}}\, e^{-\frac{x_1^2}{2}}
   \cdot 1 . \]
Per $X_2$ si ripete il calcolo scambiando i ruoli: l'espressione è simmetrica in $x_1$ e $x_2$.
\passo{La condizionata} Per definizione $f_{X_2 \mid X_1}(x_2 \mid x_1) = \frac{f_{X_1 X_2}(x_1, x_2)}{f_{X_1}(x_1)}$:
dividendo per il primo fattore resta il secondo.
\end{dimostrazione}

La densità condizionata dice due cose. La \textbf{media condizionata} $E[X_2 \mid X_1 = x_1] = \rho x_1$ è una retta:
sapendo che $X_1$ è grande, ci si aspetta che anche $X_2$ lo sia (se $\rho > 0$), ma un po' meno, perché $\abs{\rho} <
1$. La \textbf{varianza condizionata} $1 - \rho^2$ è più piccola della varianza 1 di $X_2$ e non dipende da $x_1$:
conoscere $X_1$ riduce l'incertezza su $X_2$, tanto più quanto più $\abs{\rho}$ è vicino a 1. Con $\rho = 0$ la
condizionata coincide con la marginale: conoscere $X_1$ non serve a niente, come deve essere per variabili indipendenti.

\begin{figure}[H]
\centering
\begin{tikzpicture}
\begin{groupplot}[group style={group size=2 by 1, horizontal sep=1.8cm}, height=6.2cm,
  tick label style={font=\scriptsize}, label style={font=\small}]
\nextgroupplot[width=6.6cm, xmin=-3.2, xmax=3.2, ymin=-3.2, ymax=3.2, axis equal image, xtick={-1.5,0,1.5},
  xticklabels={$-1{,}5$,$0$,$1{,}5$}, ytick={-2,-1.2,0,1.2,2}, yticklabels={$-2$,$-1{,}2$,$0$,$1{,}2$,$2$},
  xlabel=$x_1$, ylabel=$x_2$, grid=major, grid style={mspigray!20}, title={$\rho = 0{,}8$}]
  \pgfplotsinvokeforeach{1,2}{
    \addplot[mspiA, line width=0.8pt, domain=0:360, samples=73]
      ({#1*(0.9487*cos(x) + 0.3162*sin(x))}, {#1*(0.9487*cos(x) - 0.3162*sin(x))});
  }
  \addplot[mspigray, dashed, line width=0.6pt, domain=-3.2:3.2] {x};
  \addplot[black, line width=1.1pt, domain=-3.2:3.2] {0.8*x};
  \draw[mspigray, densely dotted, line width=0.6pt] (axis cs:-1.5,-3.2) -- (axis cs:-1.5,3.2)
    (axis cs:0,-3.2) -- (axis cs:0,3.2) (axis cs:1.5,-3.2) -- (axis cs:1.5,3.2);
  \addplot[only marks, mark=*, mark size=1.9pt, mspiC] coordinates {(-1.5,-1.2)};
  \addplot[only marks, mark=*, mark size=1.9pt, mspiA] coordinates {(0,0)};
  \addplot[only marks, mark=*, mark size=1.9pt, mspiB] coordinates {(1.5,1.2)};
  \node[font=\scriptsize, anchor=south east, fill=white, inner sep=1pt] at (axis cs:3.1,2.45) {$x_2 = 0{,}8\, x_1$};
\nextgroupplot[width=8.6cm, xmin=-3.6, xmax=3.6, ymin=0, ymax=0.75, xtick={-1.2,0,1.2},
  xticklabels={$-1{,}2$,$0$,$1{,}2$}, ytick=\empty, axis x line=bottom, axis y line=none, xlabel=$x_2$,
  legend style={font=\scriptsize, at={(0.5,1.03)}, anchor=south, legend columns=2, draw=none},
  every axis plot/.append style={line width=1.1pt}]
  \addplot[mspiC, domain=-3.6:3.6, samples=121] {exp(-(x+1.2)^2/(2*0.36))/sqrt(2*pi*0.36)};
  \addplot[mspiA, domain=-3.6:3.6, samples=121] {exp(-x^2/(2*0.36))/sqrt(2*pi*0.36)};
  \addplot[mspiB, domain=-3.6:3.6, samples=121] {exp(-(x-1.2)^2/(2*0.36))/sqrt(2*pi*0.36)};
  \addplot[black, dashed, domain=-3.6:3.6, samples=121] {exp(-x^2/2)/sqrt(2*pi)};
  \legend{{$x_1 = -1{,}5$}, {$x_1 = 0$}, {$x_1 = 1{,}5$}, {marginale $f_{X_2}$}}
\end{groupplot}
\end{tikzpicture}
\caption{A sinistra: le ellissi della gaussiana bivariata con $\rho = 0{,}8$ e la retta delle medie condizionate
$x_2 = 0{,}8\, x_1$ (tratteggiata la bisettrice). A destra: le densità condizionate $f_{X_2 \mid X_1}(x_2 \mid x_1)$ per
tre valori di $x_1$, tutte con varianza $1 - 0{,}8^2 = 0{,}36$ e centrate sulla retta, e la marginale $\mathcal{N}(0, 1)$,
più larga.}
\end{figure}

\subsection{Le proprietà dei vettori gaussiani}

Le slide elencano le proprietà che rendono i vettori gaussiani così comodi.
\begin{enumerate}
  \item \textbf{Bastano medie e covarianze.} Un vettore gaussiano è completamente caratterizzato dal vettore delle medie
        e dalla matrice di covarianza: nella formula della densità non compare nient'altro. Per i vettori gaussiani,
        quindi, la caratterizzazione sintetica implica quella completa; in generale è vero solo il contrario (dalla
        densità si ricavano medie e covarianze, ma non viceversa).
  \item \textbf{Chiusura rispetto alle trasformazioni lineari.} Qualsiasi trasformazione lineare di un vettore
        gaussiano è ancora un vettore gaussiano, con opportuno vettore delle medie e opportuna matrice di covarianza.
  \item \textbf{Calcoli facili.} Le medie e le covarianze dell'uscita di una trasformazione lineare si calcolano con due
        formule semplici (le vediamo nella prossima sezione); grazie alla proprietà 2, per un vettore gaussiano questo
        dà la caratterizzazione completa dell'uscita. Per vettori non gaussiani si ottiene con la stessa facilità solo la
        caratterizzazione sintetica, mentre quella completa è spesso enormemente complicata.
  \item \textbf{Sottoinsiemi.} Un qualsiasi sottoinsieme di variabili congiuntamente gaussiane è ancora un vettore
        gaussiano. In particolare ogni singola componente $X_i$ è gaussiana, con media $\mu_{X_i}$ e varianza
        $\sigma_{X_i}^2$ (per due variabili lo abbiamo appena dimostrato).
  \item \textbf{Condizionamenti.} La densità di un sottoinsieme di variabili congiuntamente gaussiane, condizionata a un
        qualsiasi sottoinsieme delle rimanenti, è ancora gaussiana, con opportune medie e covarianze condizionate (per
        due variabili: $\mathcal{N}(\rho x_1, 1 - \rho^2)$).
  \item \textbf{Incorrelazione e indipendenza coincidono} (Teorema~\ref{teo:gaussincorr}).
\end{enumerate}

% ============================================================
\section{Trasformazioni lineari di vettori gaussiani}
% ============================================================

Cominciamo dal caso più semplice: una combinazione lineare di due variabili.

\teorema[combinazione2]{Combinazione lineare di due variabili congiuntamente gaussiane}{
  Siano $X_1$ e $X_2$ congiuntamente gaussiane, con medie $\mu_1$ e $\mu_2$, varianze $\sigma_1^2$ e $\sigma_2^2$ e
  coefficiente di correlazione $\rho_{12}$, e sia $Z = a_1 X_1 + a_2 X_2$. Allora $Z \sim \mathcal{N}(\mu_Z, \sigma_Z^2)$,
  con
  \[ \mu_Z = a_1 \mu_1 + a_2 \mu_2 , \qquad \sigma_Z^2 = a_1^2 \sigma_1^2 + a_2^2 \sigma_2^2 + 2 a_1 a_2\, \sigma_1
     \sigma_2\, \rho_{12} . \]
}

\begin{dimostrazione}
\passo{La media} Per la linearità della media: $\mu_Z = E[a_1 X_1 + a_2 X_2] = a_1 \mu_1 + a_2 \mu_2$.
\passo{La varianza} Si raccolgono gli scarti dalla media e si sviluppa il quadrato del binomio:
\[ \begin{aligned}
   \sigma_Z^2 &= E[(Z - \mu_Z)^2] = E\left[ (a_1 X_1 + a_2 X_2 - a_1 \mu_1 - a_2 \mu_2)^2 \right] =
   E\left[ \left( a_1 (X_1 - \mu_1) + a_2 (X_2 - \mu_2) \right)^2 \right] \\
   &= E\left[ a_1^2 (X_1 - \mu_1)^2 + a_2^2 (X_2 - \mu_2)^2 + 2 a_1 a_2 (X_1 - \mu_1)(X_2 - \mu_2) \right] \\
   &= a_1^2 \sigma_1^2 + a_2^2 \sigma_2^2 + 2 a_1 a_2\, COV(X_1, X_2) = a_1^2 \sigma_1^2 + a_2^2 \sigma_2^2 + 2 a_1 a_2\,
   \sigma_1 \sigma_2\, \rho_{12} .
   \end{aligned} \]
Fin qui non si è usato nulla di gaussiano: sono le formule del Capitolo~9, valide per qualsiasi coppia.
\passo{La gaussianità} Che $Z$ sia gaussiana è la proprietà di chiusura; la dimostriamo alla fine della sezione, come
conseguenza del teorema seguente.
\end{dimostrazione}

Il caso generale è una trasformazione lineare qualsiasi: $m$ combinazioni lineari delle $n$ variabili, più delle
costanti.

\teorema[trasflineare]{Trasformazione lineare di un vettore gaussiano}{
  Sia $\vett{X} \sim \mathcal{N}(\vett{\mu}_{\vett{X}}, \matr{C}_{\vett{X}})$ un vettore gaussiano $n \times 1$, e sia
  \[ \underbrace{\vett{Z}}_{m \times 1} = \underbrace{\matr{A}}_{m \times n}\, \underbrace{\vett{X}}_{n \times 1} +
     \underbrace{\vett{b}}_{m \times 1} , \]
  con $\matr{A}$ e $\vett{b}$ costanti. Allora $\vett{Z} \sim \mathcal{N}(\vett{\mu}_{\vett{Z}}, \matr{C}_{\vett{Z}})$,
  con
  \[ \underbrace{\vett{\mu}_{\vett{Z}}}_{m \times 1} = \underbrace{\matr{A}}_{m \times n}\,
     \underbrace{\vett{\mu}_{\vett{X}}}_{n \times 1} + \underbrace{\vett{b}}_{m \times 1} , \qquad
     \underbrace{\matr{C}_{\vett{Z}}}_{m \times m} = \underbrace{\matr{A}}_{m \times n}\,
     \underbrace{\matr{C}_{\vett{X}}}_{n \times n}\, \underbrace{\matr{A}^T}_{n \times m} . \]
  Le formule per $\vett{\mu}_{\vett{Z}}$ e $\matr{C}_{\vett{Z}}$ valgono per qualsiasi vettore aleatorio $\vett{X}$,
  gaussiano o no.
}

\begin{dimostrazione}
\passo{La media} Ogni componente di $\vett{Z}$ è una combinazione lineare delle $X_j$ più una costante,
$Z_i = \sum_j a_{ij} X_j + b_i$. Per la linearità della media $E[Z_i] = \sum_j a_{ij} \mu_{X_j} + b_i$, cioè
$\vett{\mu}_{\vett{Z}} = \matr{A}\vett{\mu}_{\vett{X}} + \vett{b}$.
\passo{Lo scarto dalla media} La costante $\vett{b}$ si cancella:
$\vett{Z} - \vett{\mu}_{\vett{Z}} = \matr{A}\vett{X} + \vett{b} - \matr{A}\vett{\mu}_{\vett{X}} - \vett{b} =
\matr{A}(\vett{X} - \vett{\mu}_{\vett{X}})$.
\passo{La covarianza} Per il trasposto di un prodotto,
$\left( \matr{A}(\vett{X} - \vett{\mu}_{\vett{X}}) \right)^T = (\vett{X} - \vett{\mu}_{\vett{X}})^T \matr{A}^T$, e le
matrici costanti escono dalla media (una a sinistra e una a destra):
\[ \matr{C}_{\vett{Z}} = E\left[ (\vett{Z} - \vett{\mu}_{\vett{Z}})(\vett{Z} - \vett{\mu}_{\vett{Z}})^T \right] =
   E\left[ \matr{A} (\vett{X} - \vett{\mu}_{\vett{X}}) (\vett{X} - \vett{\mu}_{\vett{X}})^T \matr{A}^T \right] =
   \matr{A}\, E\left[ (\vett{X} - \vett{\mu}_{\vett{X}}) (\vett{X} - \vett{\mu}_{\vett{X}})^T \right] \matr{A}^T =
   \matr{A} \matr{C}_{\vett{X}} \matr{A}^T . \]
\passo{La gaussianità, con $m = n$ e $\matr{A}$ invertibile} Si ragiona come per $Y = aX + b$ nel Capitolo~12. Ogni
$\vett{z}$ viene da un solo $\vett{x} = \matr{A}^{-1}(\vett{z} - \vett{b})$. Una scatolina attorno a $\vett{x}$ viene
trasformata in una scatolina attorno a $\vett{z}$ con volume $\abs*{\det\matr{A}}$ volte più grande (richiamo), che
contiene la stessa probabilità; quindi la densità, che è una probabilità per unità di volume, si divide per
$\abs*{\det\matr{A}}$:
\[ f_{\vett{Z}}(\vett{z}) = \frac{1}{\abs*{\det \matr{A}}}\, f_{\vett{X}}\left( \matr{A}^{-1}(\vett{z} - \vett{b})
   \right) . \]
Nell'esponente compare $\vett{x} - \vett{\mu}_{\vett{X}} = \matr{A}^{-1}(\vett{z} - \vett{b}) -
\matr{A}^{-1}\matr{A}\vett{\mu}_{\vett{X}} = \matr{A}^{-1}(\vett{z} - \vett{\mu}_{\vett{Z}})$, e la forma quadratica
diventa
\[ (\vett{z} - \vett{\mu}_{\vett{Z}})^T (\matr{A}^{-1})^T \matr{C}_{\vett{X}}^{-1} \matr{A}^{-1} (\vett{z} -
   \vett{\mu}_{\vett{Z}}) = (\vett{z} - \vett{\mu}_{\vett{Z}})^T \left( \matr{A} \matr{C}_{\vett{X}} \matr{A}^T
   \right)^{-1} (\vett{z} - \vett{\mu}_{\vett{Z}}) = (\vett{z} - \vett{\mu}_{\vett{Z}})^T \matr{C}_{\vett{Z}}^{-1}
   (\vett{z} - \vett{\mu}_{\vett{Z}}) , \]
perché l'inversa di un prodotto è il prodotto delle inverse in ordine inverso:
$(\matr{A}\matr{C}_{\vett{X}}\matr{A}^T)^{-1} = (\matr{A}^T)^{-1} \matr{C}_{\vett{X}}^{-1} \matr{A}^{-1}$. Infine
$\det \matr{C}_{\vett{Z}} = \det\matr{A} \cdot \det\matr{C}_{\vett{X}} \cdot \det\matr{A}^T = (\det\matr{A})^2
\det\matr{C}_{\vett{X}}$, quindi $\abs*{\det\matr{A}} \sqrt{(2\pi)^n \det\matr{C}_{\vett{X}}} = \sqrt{(2\pi)^n
\det\matr{C}_{\vett{Z}}}$. Mettendo insieme, $f_{\vett{Z}}$ ha proprio la forma gaussiana, con media
$\vett{\mu}_{\vett{Z}}$ e covarianza $\matr{C}_{\vett{Z}}$.
\passo{Meno uscite che ingressi ($m < n$)} Si aggiungono ad $\matr{A}$ delle righe in modo da renderla quadrata e
invertibile: il vettore allungato è gaussiano per il passo precedente, e $\vett{Z}$, che ne è un sottoinsieme, è
gaussiano per la proprietà dei sottoinsiemi. (Se invece $m > n$, o $\matr{A}$ non è invertibile, si trova
$\det\matr{C}_{\vett{Z}} = 0$: il vettore $\vett{Z}$ sta tutto su una retta o su un piano e una densità non si può
scrivere. È un caso «degenere» che non ci servirà.)
\end{dimostrazione}

Ora possiamo completare la dimostrazione del Teorema~\ref{teo:combinazione2}. La combinazione $Z = a_1 X_1 + a_2 X_2$ è
il caso $m = 1$, $n = 2$, $\matr{A} = [a_1, a_2]$, $\vett{b} = 0$, e la formula della covarianza ridà la varianza (è la
forma quadratica della sezione precedente):
\[ \sigma_Z^2 = [a_1, a_2] \begin{bmatrix} \sigma_1^2 & \rho_{12}\sigma_1\sigma_2 \\ \rho_{12}\sigma_1\sigma_2 &
   \sigma_2^2 \end{bmatrix} \begin{bmatrix} a_1 \\ a_2 \end{bmatrix} = a_1^2 \sigma_1^2 + a_2^2 \sigma_2^2 + 2 a_1
   a_2\, \sigma_1 \sigma_2\, \rho_{12} . \]
Per la gaussianità, se $a_1 \neq 0$ si aggiunge la riga $[0, 1]$, che rende la matrice quadrata e invertibile:
\[ \begin{bmatrix} Z \\ X_2 \end{bmatrix} = \begin{bmatrix} a_1 & a_2 \\ 0 & 1 \end{bmatrix} \begin{bmatrix} X_1 \\
   X_2 \end{bmatrix} , \qquad \det \begin{bmatrix} a_1 & a_2 \\ 0 & 1 \end{bmatrix} = a_1 \cdot 1 - a_2 \cdot 0 = a_1
   \neq 0 . \]
Il vettore $[Z, X_2]^T$ è gaussiano (passo 4) e quindi $Z$, che ne è una componente, è gaussiana (lo garantisce il
Teorema~\ref{teo:gaussmarginale}, dopo aver standardizzato). Se $a_1 = 0$, $Z = a_2 X_2$ è una gaussiana moltiplicata
per una costante, che è gaussiana per il Capitolo~15.

\begin{esempio}{Somma e differenza di due rumori indipendenti}
Due campioni di rumore $X_1$ e $X_2$ sono gaussiani indipendenti, entrambi $\mathcal{N}(0, \sigma^2)$: il vettore
$\vett{X}$ è gaussiano con $\vett{\mu}_{\vett{X}} = \vett{0}$ e $\matr{C}_{\vett{X}} = \sigma^2 \matr{I}$. Calcoliamo
somma e differenza:
\[ \vett{Z} = \begin{bmatrix} X_1 + X_2 \\ X_1 - X_2 \end{bmatrix} = \underbrace{\begin{bmatrix} 1 & 1 \\ 1 & -1
   \end{bmatrix}}_{\matr{A}} \begin{bmatrix} X_1 \\ X_2 \end{bmatrix} , \qquad
   \matr{C}_{\vett{Z}} = \matr{A}\, (\sigma^2 \matr{I})\, \matr{A}^T = \sigma^2 \matr{A}\matr{A}^T = \sigma^2
   \begin{bmatrix} 1 \cdot 1 + 1 \cdot 1 & 1 \cdot 1 + 1 \cdot (-1) \\ 1 \cdot 1 + (-1) \cdot 1 & 1 \cdot 1 + (-1)(-1)
   \end{bmatrix} = \begin{bmatrix} 2\sigma^2 & 0 \\ 0 & 2\sigma^2 \end{bmatrix} . \]
La somma e la differenza sono gaussiane con varianza $2\sigma^2$ e sono \textbf{incorrelate}; essendo congiuntamente
gaussiane (Teorema~\ref{teo:trasflineare}), sono anche \textbf{indipendenti} (Teorema~\ref{teo:gaussincorr}). Conoscere
la somma dei due rumori non dice nulla sulla loro differenza! Geometricamente, $\matr{A}$ non deforma le figure: le
ribalta come uno specchio (attorno a una retta inclinata) e le ingrandisce di $\sqrt{2}$. Una campana rotonda,
ribaltata e ingrandita, resta rotonda. Con altre distribuzioni non succede: se $X_1$ e $X_2$
sono uniformi su $[0, 1]$, somma e differenza sono ancora incorrelate (la covarianza è $\sigma_{X_1}^2 - \sigma_{X_2}^2
= 0$), ma non indipendenti: se la somma è vicina a 0, entrambe le variabili sono vicine a 0, e anche la differenza lo è.
\end{esempio}

\begin{esempio}{Come si costruisce un vettore gaussiano correlato}
Un calcolatore sa generare numeri gaussiani standard indipendenti, $W_1$ e $W_2$. Come si ottiene una coppia gaussiana
con varianze unitarie e correlazione $\rho$ assegnata? Si prende
\[ \begin{bmatrix} X_1 \\ X_2 \end{bmatrix} = \underbrace{\begin{bmatrix} 1 & 0 \\ \rho & \sqrt{1 - \rho^2}
   \end{bmatrix}}_{\matr{A}} \begin{bmatrix} W_1 \\ W_2 \end{bmatrix} , \qquad\text{cioè}\qquad X_1 = W_1 , \quad X_2 =
   \rho\, W_1 + \sqrt{1 - \rho^2}\, W_2 . \]
Il vettore $\vett{X}$ è gaussiano (trasformazione lineare di un vettore gaussiano), con media nulla e
\[ \matr{C}_{\vett{X}} = \matr{A}\, \matr{I}\, \matr{A}^T = \begin{bmatrix} 1 & 0 \\ \rho & \sqrt{1 - \rho^2}
   \end{bmatrix} \begin{bmatrix} 1 & \rho \\ 0 & \sqrt{1 - \rho^2} \end{bmatrix} = \begin{bmatrix} 1 & \rho \\ \rho &
   \rho^2 + 1 - \rho^2 \end{bmatrix} = \begin{bmatrix} 1 & \rho \\ \rho & 1 \end{bmatrix} , \]
proprio la matrice voluta. È così che sono stati generati i punti delle figure di questo capitolo. Si ritrova anche la
condizionata del Teorema~\ref{teo:gaussmarginale}: sapendo che $X_1 = W_1 = x_1$, si ha $X_2 = \rho x_1 + \sqrt{1 -
\rho^2}\, W_2$, una gaussiana di media $\rho x_1$ e varianza $1 - \rho^2$ ($W_2$ è indipendente da $W_1$).
\end{esempio}

\begin{esempio}{La media di $n$ campioni di rumore}
Siano $X_1, \ldots, X_n$ gaussiane indipendenti, tutte $\mathcal{N}(\mu, \sigma^2)$ (per esempio $n$ misure della stessa
grandezza, disturbate da rumore). La loro media aritmetica è una trasformazione lineare con $m = 1$:
\[ \bar{X} = \frac{X_1 + \cdots + X_n}{n} = \underbrace{\left[ \tfrac1n, \tfrac1n, \ldots, \tfrac1n
   \right]}_{\matr{A}\ (1 \times n)} \vett{X} . \]
Quindi $\bar{X}$ è gaussiana, con media $\matr{A}\vett{\mu}_{\vett{X}} = n \cdot \frac{\mu}{n} = \mu$ e varianza
\[ \matr{A}\, (\sigma^2 \matr{I})\, \matr{A}^T = \sigma^2 \sum_{i=1}^{n} \frac{1}{n^2} = \frac{\sigma^2}{n} . \]
Mediando $n$ misure la deviazione standard del rumore si divide per $\sqrt{n}$: con 100 misure l'errore tipico diventa
dieci volte più piccolo. Ritroveremo questa media, la \emph{media campione}, all'inizio del terzo pacchetto di slide.
\end{esempio}

\gbox{In sintesi}{mspiA!80!black}{
\begin{itemize}
  \item Un \textbf{vettore aleatorio} $\vett{X} = [X_1, \ldots, X_n]^T$ ha vettore delle medie $\vett{\mu}_{\vett{X}} =
        E[\vett{X}]$, matrice di correlazione $\matr{R}_{\vett{X}} = E[\vett{X}\vett{X}^T]$ e matrice di covarianza
        $\matr{C}_{\vett{X}} = E[(\vett{X} - \vett{\mu}_{\vett{X}})(\vett{X} - \vett{\mu}_{\vett{X}})^T] =
        \matr{R}_{\vett{X}} - \vett{\mu}_{\vett{X}}\vett{\mu}_{\vett{X}}^T$, entrambe simmetriche; sulla diagonale di
        $\matr{C}_{\vett{X}}$ ci sono le varianze. $(\vett{\mu}_{\vett{X}}, \matr{C}_{\vett{X}})$ è la caratterizzazione
        sintetica.
  \item $\matr{C}_{\vett{X}}$ è \textbf{definita non negativa}: $\vett{q}^T \matr{C}_{\vett{X}} \vett{q} =
        \mathrm{VAR}(\vett{q}^T \vett{X}) \geq 0$.
  \item \textbf{Vettore gaussiano}, $\vett{X} \sim \mathcal{N}(\vett{\mu}_{\vett{X}}, \matr{C}_{\vett{X}})$:
        \[ f_{\vett{X}}(\vett{x}) = \frac{1}{\sqrt{(2\pi)^n \det \matr{C}_{\vett{X}}}}\; e^{-\frac12 (\vett{x} -
           \vett{\mu}_{\vett{X}})^T \matr{C}_{\vett{X}}^{-1} (\vett{x} - \vett{\mu}_{\vett{X}})} ; \]
        con componenti indipendenti $\matr{C}_{\vett{X}}$ è diagonale. Le curve di uguale densità sono ellissi.
  \item Per le congiuntamente gaussiane \textbf{incorrelate $\Leftrightarrow$ indipendenti}; marginali e condizionate
        sono gaussiane (bivariata standard: $X_2 \mid X_1 = x_1 \sim \mathcal{N}(\rho x_1, 1 - \rho^2)$).
  \item \textbf{Trasformazioni lineari}: $\vett{Z} = \matr{A}\vett{X} + \vett{b}$ ha $\vett{\mu}_{\vett{Z}} =
        \matr{A}\vett{\mu}_{\vett{X}} + \vett{b}$ e $\matr{C}_{\vett{Z}} = \matr{A}\matr{C}_{\vett{X}}\matr{A}^T$, ed è
        gaussiano se $\vett{X}$ lo è. Per due variabili: $\sigma_Z^2 = a_1^2\sigma_1^2 + a_2^2\sigma_2^2 +
        2a_1a_2\sigma_1\sigma_2\rho_{12}$.
\end{itemize}
}
